\subsection{Posterior modelling}
\label{sec:posterior_modelling}

To address the limited identifiability of a single global discrepancy factor and to make better use of the probabilistic input models in Table~2, Bayesian updating was reformulated at the level of the basic variables. The prior vector $\boldsymbol{X}$ contains the random variables listed in Table~2. For the fib Bulletin~34 model these include $R^{-1}_{ACC,0}$, $k_t$, $\varepsilon_t$, $t_c$, $b_c$, $[CO_2]_{ppm}$, $RH$, $t_w$, $b_w$, $p_{SR}$ and $\theta_{xc}$. For the Malami model these include $W/C$, $t_c$, $b_c$, $[CO_2]_{ppm}$, $RH$, $t_w$, $b_w$, $p_{SR}$, $T$ and $\theta_{xc}$, while the mixture-composition quantities without direct inspection information are kept at the deterministic values reported in Table~2. Thus, the default probabilistic models in Table~2 are used as priors rather than treating only $\theta_{xc}$ as uncertain.

For a given sample of $\boldsymbol{X}$, the carbonation depth is evaluated as
\begin{equation}
    x_c(t,\boldsymbol{X}) = \theta_{xc} A(\boldsymbol{X})t^{0.5-n(\boldsymbol{X})},
\end{equation}
where $A(\boldsymbol{X})$ is obtained from the corresponding carbonation model and $n(\boldsymbol{X})=(p_{SR}t_w)^{b_w}/2$. The inspection result at $t=42$ years is introduced through a Gaussian likelihood,
\begin{equation}
    p(\bar{x}_{c,obs}\mid\boldsymbol{X}) =
    \mathcal{N}\left(\bar{x}_{c,obs};x_c(42,\boldsymbol{X}),\sigma_{obs}^2/N_{eff}+\sigma_\varepsilon^2\right),
\end{equation}
where $\bar{x}_{c,obs}=14.78$ mm and $\sigma_{obs}=5.3$ mm describe the carbonation-depth statistics from the inspection. $N_{eff}$ denotes the effective number of independent observations. In the absence of the full spatial measurement vector, $N_{eff}=1$ is adopted to avoid over-conditioning the posterior on a single time-point summary. If the individual measurements are used, the likelihood is replaced by the product of the corresponding Gaussian observation densities.

The posterior distribution is proportional to the product of the likelihood and the Table~2 priors,
\begin{equation}
    p(\boldsymbol{X}\mid\bar{x}_{c,obs}) \propto
    p(\bar{x}_{c,obs}\mid\boldsymbol{X})\prod_j p(X_j).
\end{equation}
Sampling is performed using a random-walk Metropolis Markov chain Monte Carlo algorithm in transformed standard space. Lognormal variables are sampled through their logarithmic normal representation, Gaussian variables through standardisation, and bounded uniform variables through a logit transformation. At each iteration a candidate vector is generated by adding a zero-mean Gaussian perturbation to the current transformed state. The candidate is accepted with probability
\begin{equation}
    \alpha = \min\left[1,\frac{p(\boldsymbol{X}^{*}\mid\bar{x}_{c,obs})}{p(\boldsymbol{X}^{(k)}\mid\bar{x}_{c,obs})}\right].
\end{equation}
The first samples are discarded as burn-in and the retained chain is thinned before estimating posterior means, standard deviations and predictive carbonation-depth trajectories.

This updating format does not force the model prediction to match the 42-year observation through one corrective variable only. Instead, the inspection information is allowed to update the most relevant uncertain material, exposure and model-discrepancy variables, while the prior distributions prevent unrealistic compensation among poorly identifiable parameters. The resulting posterior therefore represents a case-specific calibration of the probabilistic input model rather than a deterministic fit to a single inspection point.
